\section{Related Literature}
\label{sec:literature}

\paragraph{Structural estimation in auctions.}
The structural auction literature recovers latent primitives from
observed bids \citep{guerre2000optimal, athey2001information}.
\citet{porter1993detection} and \citet{pesendorfer2000study} use
structural methods to detect bid rigging in procurement;
\citet{krasnokutskaya2011bid} shows that unobserved auction
heterogeneity biases cost estimates when ignored.
We contribute a \emph{mixture model with observed labels}: the FL
classification partitions bids into genuine and cover components
without requiring an EM algorithm or latent-type assumptions.

\paragraph{Bid-coordination tests.}
\citet{bajari2003deciding} develop exchangeability and
conditional-independence tests to distinguish competitive from
collusive bidding; applications require specifying bidder groups
\emph{a priori} \citep{conley2016detecting}. FL status provides a
data-driven partition. The structural model adds an insight absent
from the original Bajari--Ye framework: the \emph{direction} of
the exchangeability violation depends on the cover-bidding
regime, and the tender-FE reversal is a novel prediction that we
find consistent with the data in Section~\ref{sec:spec_tests}.

\paragraph{Cartel screens and the dispersion assumption.}
Statistical screens typically flag high bid dispersion
\citep{abrantesmetz2006a, imhof2019detecting, huber2019machine};
\citet{chassang2022robust} design screens robust to strategic
manipulation. These approaches implicitly assume Regime~1
(complementary cover bidding). Our structural estimation selects
Regime~2 (coordinated) in this setting, providing one
interpretation of the performance gap between participation-based
and variance-based screens.

\paragraph{Cover bidding in practice.}
\citet{asker2010leniency} documents Regime~1 patterns in a
philatelic stamp cartel; \citet{marshall2012economics} note that
sophisticated cartels instruct cover bidders to bid ``just above''
the winner---a Regime~2 pattern. We provide the first structural
test distinguishing these regimes.

\paragraph{Procurement design.}
Minimum-bidder requirements are common worldwide
\citep{szymanski2006impact, decarolis2018buyers}. Our model shows
that Brazil's three-bidder rule for convite creates a corner solution
(Proposition~\ref{prop:optimal_m}) that forces cover-bidder
deployment. A companion paper \citep{genicolomartins2026screening}
validates the FL screen; the present paper develops the structural
mechanism.
